\documentclass[11pt]{exam}

%%%%%%%%%%%%%%%%%%%%%%%%%% Preamble %%%%%%%%%%%%%%%%%%%%%%%%%% 
\input{01. Admin/1. Packages}
\input{01. Admin/2. WEB Course Info}
\input{01. Admin/3. WEB Outcomes}
\input{01. Admin/4. Commands}
\input{01. Admin/8. Fonts & Display}

%%%%%%%%%%%%% Exam Information  %%%%%%%%%%%%%%%%%%%%%%%%%% 

\newcommand{\exam}{Comprehensive Final}
\newcommand{\TimeLimit}{2 hours}
\newcommand{\version}{SP26A}

\runningheader{}{}{\version}

\begin{document}

%%%% Show/Don't Show Points %%%% 
\addpoints
% \pointformat{}
% \nopointsmargin

%%%% Show/Don't Show Solutions %%%% 
%\printanswers
\noprintanswers

%%%%%%%%%%%%%%%%%%%%%%%%%% Scratch Paper %%%%%%%%%%%%%%%%%%%%%%%%%%
\blankpage

%%%%%%%%%%%%%%%% Cover Page %%%%%%%%%%%%%%%% 
\thispagestyle{empty}
\input{01. Admin/6. F2F Final CoverPage}

\newpage

%%%% Questions %%%%


\begin{questions}

\question[10] Evaluate and simplify the difference quotient, $\dfrac{f(x+h)-f(x)}{h}$, for the function. \\ \\ 
\(f(x) = 5 \, x^{2} + 8\)
    \begin{solution}
         \(f(x+h) = 5(x+h)^2 + 8\) 
         \(= 5(x^2 + 2xh + h^2) + 8\) 
         \(= 5 \, h^{2} + 10 \, h x + 5 \, x^{2} + 8\) 
         \(\displaystyle{\frac{f(x+h)-f(x)}{h} = \frac{ (5 \, h^{2} + 10 \, h x + 5 \, x^{2} + 8) - (5 \, x^{2} + 8) }{h}}\) 
         \(=\displaystyle{ \frac{ 5 \, h^{2} + 10 \, h x }{h}}\) 
         \(= \displaystyle{ \frac{ h(5 \, h + 10 \, x) }{h}}\) 
         \(\boxed{ 5 \, h + 10 \, x }\) 
    \end{solution}
\vspace{\stretch{1}}\answerline

\newpage

\question[8] Use the table to identify the transformations of the toolkit function described by the equation. Circle the option that applies and fill in the blanks as appropriate to describe the transformations on the given function.  If one does not apply, you may leave it blank. Then sketch the transformed graph and state the resulting domain and range.\\

\(f(x) = \log_{3}(x + 6) - 1\)\\

\renewcommand{\arraystretch}{2.5}

	\begin{tabular}{|l|l|}
		\hline
		\textbf{Horizontal Transformations}               & \textbf{Vertical Transformations}              \\ \hline
		Reflection: YES    or     NO       & Reflection:       YES    or     NO    \\ \hline
		Dilation: \underline{\hspace{2cm}} times as wide           & Dilation: \underline{\hspace{2cm}} times as tall        \\ \hline
		Translation: \underline{\hspace{2cm}} units LEFT or RIGHT & Translation: \underline{\hspace{2cm}} units UP or DOWN \\ \hline
	\end{tabular}

\vspace{\stretch{1}}
\begin{center}
    \begin{tikzpicture}[scale=0.5,>=triangle 45]
      \draw [step=1cm, style={black!60}] (-10,-10) grid (10,10);
      \draw [<->, very thick] (-10.5,0) -- (10.5,0); 
      \draw [<->, very thick] (0,-10.5) -- (0,10.5);
    \end{tikzpicture}
\end{center}
\vspace{\stretch{1}}

Domain: \fillin[][2in] \hspace{\stretch{1}} Range: \fillin[][2in]

\begin{solution}
    \begin{itemize}
\item
 \(\textbf{Horizontal: } \text{Reflection: NO, Dilation: 1, Shift: 6 units LEFT.}\) 

\item
 \(\textbf{Vertical: } \text{Reflection: NO, Dilation: 1, Shift: 1 units DOWN.}\) 

\end{itemize}


 \[\begin{tikzpicture}[scale=0.39,>=triangle 45] \draw[step=1cm, black!60] (-10,-10) grid (10,10); \draw[very thick, <->] (-10.5,0) -- (10.5,0); \draw[very thick, <->] (0,-10.5) -- (0,10.5); \clip (-10.5,-10.5) rectangle (10.5,10.5);\draw[dashed, red, thick] (-6, -10.5) -- (-6, 10.5);\draw[blue, very thick, samples=100, domain=-5.99000000000000:10.5000000000000, smooth] plot (\x, {(ln(\x - (-6))/ln(3)) + (-1)});\end{tikzpicture}\] 
\end{solution}

\newpage

\question[5] Solve each of the following equations. Give exact solutions.\\
\begin{parts}
\part
\(\log_{5}(x + 2) + \log_{5}(x + 6) = 1\)
\begin{solution}
    \[\begin{aligned} \log_{5}((x + 2)(x + 6)) &= 1 \\ 5 &= (x + 2)(x + 6) \\ 5 &= x^{2} + 8 \, x + 12 \\ 0 &= x^{2} + 8 \, x + 7 \\ 0 &= (x + 1)(x + 7) \end{aligned}\]

 \(x = -1\) is a valid solution, while \(x = -7\) causes a non-positive argument in the original logarithms. 

solutions: \(\boxed{ x = -1 }\)

extraneous solutions: \(\boxed{ x = -7 }\)
\end{solution}
    \vspace{\stretch{1}}\answerline

\part
\(9^{x + 3} \cdot 9^{-3 \, x - 2} = 9^{-3}\)
\begin{solution}
    \[\begin{aligned} 9^{x + 3-3 \, x - 2} &= 9^{-3} \\ x + 3-3 \, x - 2 &= -3 \\ -2 \, x + 1 &= -3 \\ -2 \, x &= -4 \\ x &= 2 \end{aligned}\]

solutions: \(\boxed{ 2 }\)

extraneous solutions: \(\boxed{ \text{none} }\)
\end{solution}
    
    \vspace{\stretch{1}}\answerline

\end{parts}

\newpage

\question[5] Use logarithm identities to determine the exact value of the expression, given \(\log_b(m)=7\), \(\log_b(n-1)=10\), and \(\log_b(p)=7\).\\ \\

\(\log_b\left(\dfrac{m^{3}\sqrt{n-1}}{p^{2}}\right)\)
\begin{solution}
    \(3(7) + \frac{1}{2}(10) - 2(7)\)
    \(21 + 5 - 14\)
    \(\boxed{12}\)
\end{solution}

\vspace{\stretch{1}}\answerline

\newpage

\question[5] Right triangle $ABC$ has the following measurements: $c=20$ in and $A=44.9^\circ$.\\
\begin{parts}
    \part Sketch the triangle.
        \vspace{\stretch{1}}
    \part Solve $\triangle ABC$.\\
           \hspace*{2.5in}
        \renewcommand{\arraystretch}{2.5}
       \begin{tabular}{|p{4cm}|p{4cm}|}\hline
         $A=$  & $a=$ \\\hline
        $B=$  & $b=$ \\\hline
        $C=$  & $c=$ \\\hline
      \end{tabular}
        \vspace{\stretch{1}}   
\end{parts}

\begin{solution}
    \(B = 90^\circ - 44.9^\circ = 45.1^\circ\)

\(\sin(44.9^\circ) = \frac{a}{20} \implies a = 20\sin(44.9^\circ) \approx 14.12\)

\(\cos(44.9^\circ) = \frac{b}{20} \implies b = 20\cos(44.9^\circ) \approx 14.17\)

\(\boxed{\begin{aligned} A &= 44.9^\circ & a &\approx 14.12 \\ B &= 45.1^\circ & b &\approx 14.17 \\ C &= 90^\circ & c &\approx 20.00 \end{aligned}}\)

\end{solution}

\newpage

\question[5]
Verify the given identity.\\ \\
\(\tan(2t) = \dfrac{2\sin(t)\cos(t)}{\cos^2(t) - \sin^2(t)}\)
    \begin{solution}
    
        \(\frac{2\sin(t)\cos(t)}{\cos^2(t) - \sin^2(t)}\)
        \(= \frac{\sin(2t)}{\cos(2t)}\)
        \(= \boxed{\tan(2t)}\)
    \end{solution}
\vspace{\stretch{1}}

\newpage

\question[5]
Find all solutions on the interval $[0^\circ,360^\circ)$. Round approximate solutions to the nearest tenth of a degree as appropriate.\\ \\
\(2 \cos \theta = -\dfrac{7}{4}\)
    \begin{solution}
        \(\cos \theta = -\frac{7}{8}\)
    
        \(\theta = \cos^{-1}\left(-\frac{7}{8}\right)\)
        \(\theta \approx 151.045^\circ \dots\)
        \(360^\circ - 151.045^\circ \approx 208.955^\circ\)
        \(\boxed{\theta = 151.0^\circ, 209.0^\circ}\)
    \end{solution}
\vspace{\stretch{1}}
\answerline

\question[5] Solve for \textit{\textbf{all possible}} triangles. Round approximate values to the nearest tenth as appropriate. \emph{(Note: If only one possible triangle exists, you may leave the other table blank.)}

$\triangle{ABC}$: \quad $A=28^\circ$ \qquad $a=8.13$ in \qquad $b=13.21$ in
\begin{solution}
    \(\text{Triangle 1:}\)
    \(\sin(B_1) = \frac{b\sin(A)}{a}\)
    \(B_1 = \sin^{-1}\left(\frac{13.21\sin(28^\circ)}{8.13}\right) \approx 49.7^\circ\)
    \(C_1 = 180^\circ - 28^\circ - 49.7^\circ = 102.3^\circ\)
    \(c_1 = \frac{a\sin(C_1)}{\sin(A)}\)
    \(c_1 = \frac{8.13\sin(102.3^\circ)}{\sin(28^\circ)} \approx 16.92\)
    \(\text{Triangle 2:}\)
    \(B_2 = 180^\circ - B_1 \approx 130.3^\circ\)
    \(C_2 = 180^\circ - 28^\circ - 130.3^\circ = 21.7^\circ\)
    \(c_2 = \frac{a\sin(C_2)}{\sin(A)}\)
    \(c_2 = \frac{8.13\sin(21.7^\circ)}{\sin(28^\circ)} \approx 6.41\)
\end{solution}

\begin{multicols}{2}

    \renewcommand{\arraystretch}{2.5}
\begin{tabular}{|p{3cm}|p{3cm}|}\hline
   $A=$  & $a=$ \\\hline
   $B=$  & $b=$ \\\hline
   $C=$  & $c=$ \\\hline
\end{tabular}

\columnbreak

\renewcommand{\arraystretch}{2.5}
\begin{tabular}{|p{3cm}|p{3cm}|}\hline
   $A=$  & $a=$ \\\hline
   $B=$  & $b=$ \\\hline
   $C=$  & $c=$ \\\hline
\end{tabular}

\end{multicols}
\vspace{\stretch{1}}

\newpage

\question[3]
A city is developing a new public park in a triangular lot left over from a new highway-interchange project. The three sides of the lot are bordered by existing roads and measure approximately \(570\) feet on the first street, \(493\) feet on the second street, and \(553\) feet along the highway. How many square feet of sod need to be ordered to cover the entire lot?
\begin{solution}
    \(s = \frac{570 + 493 + 553}{2} = 808\)
    \(A = \sqrt{s(s-570)(s-493)(s-553)}\)
    \(A = \sqrt{808(238)(315)(255)}\)
    \(A \approx 124,285.23\text{ ft}^2\)
\end{solution}

\vspace{\stretch{1}}
\answerline

\question[3]
Eliminate the parameter of the pair of parametric equations.\\ \\
\(x = 3 \, \cos\left(t\right) \qquad y = 15 \, \cos\left(t\right) + 3\)
    \begin{solution}
        \(\cos\left(t\right) = \frac{1}{3} \, x\)
        \(y = 15 \, \left(\frac{1}{3} \, x\right) + 3\)
        \(y = 5 \, x + 3\)
    \end{solution}
\vspace{\stretch{1}}
\answerline

\newpage

\question[6]{Determine each of the following limits. If the limit does not exist, write \emph{DNE}.}
\begin{parts}
   \part 
    \(\displaystyle \lim_{x \to 8} \sqrt{x^{2} - 10 \, x + 25}\)
\begin{solution}
    \(\displaystyle \lim_{x \to 8} \sqrt{x^{2} - 10 \, x + 25} = \sqrt{9} = 3\)
\end{solution}
\vspace{\stretch{1}}
\answerline
   \part 
    \(\displaystyle \lim_{x \to (-2)} \frac{x^{2} + 3 \, x + 2}{x^{2} - 4}\)
\begin{solution}
    \(\displaystyle \lim_{x \to (-2)} \frac{x^{2} + 3 \, x + 2}{x^{2} - 4} = \displaystyle \lim_{x \to (-2)} \frac{(x + 2)(x + 1)}{(x + 2)(x - 2)} = \displaystyle \lim_{x \to (-2)} \frac{x + 1}{x - 2} = \frac{1}{4}\)
\end{solution}
\vspace{\stretch{1}}
\answerline
   \part 
    \(\displaystyle \lim_{x \to (-4)} \frac{x^{2} + 2 \, x - 4}{x + 5}\)
\begin{solution}
    \(\displaystyle \lim_{x \to (-4)} \frac{x^{2} + 2 \, x - 4}{x + 5} = \frac{4}{1} = 4\)
\end{solution}
\vspace{\stretch{1}}
\answerline
\end{parts}

\calcfoot        

\newpage

% \blankpage
\includepdf[]{Formula Sheet/Formulas.pdf}
\blankpage
\blankpage

\nocalc{\version}

\setcounter{question}{0}

\question[3] Suppose a \textbf{rational} function has the attributes in the table given below.  Using only the holes, asymptotes, and intercepts listed, along with as many of the helpful points as required, sketch the graph of the function. \textit{\textbf{Be careful not to include any intercepts, holes, or asymptotes other than the ones listed.}}\\

\begin{multicols}{2}
\renewcommand{\arraystretch}{2}
\,\\
\begin{tabular}{ |r|c c| } 
 \hline
 Hole: & $\left(-8,-8.9\right)$ & \\ \hline
 Asymptotes: & $x=(-5)$ & $x=4$ \\
  & $y=(-4)$& \\ \hline
 Intercepts: & $(8,0)$ & $(-3,0)$ \\
 & $(0, -4.8)$ & \\\hline
 Helpful Points: & $(-9,-7.9)$ & $(-4,6)$ \\ 
 & $\left(6.5,2\right)$ & \\
 \hline
\end{tabular}

\columnbreak

    \begin{tikzpicture}[scale=0.4,>=triangle 45]
      \draw [step=1cm, style={black!60}] (-10,-10) grid (10,10);
      \draw [<->, very thick] (-10.5,0) -- (10.5,0); 
      \draw [<->, very thick] (0,-10.5) -- (0,10.5);
    \end{tikzpicture}
    \end{multicols}

\begin{solution}
     \textbf{Instructor Note:} The generated function is: \[f(x) = \frac{-4(x - 8)(x + 3)(x + 8)}{(x + 5)(x - 4)(x + 8)}\] 

 \[\begin{tikzpicture}[scale=0.39,>=triangle 45] \draw[step=1cm, black!60] (-10,-10) grid (10,10); \draw[very thick, <->] (-10.5,0) -- (10.5,0); \draw[very thick, <->] (0,-10.5) -- (0,10.5); \clip (-10.5,-10.5) rectangle (10.5,10.5);\draw[dashed, blue, very thick, <->] (-5, -10.5) -- (-5, 10.5); \draw[dashed, blue, very thick, <->] (4, -10.5) -- (4, 10.5); \draw[dashed, blue, very thick, <->] (-10.5, -4) -- (10.5, -4); \draw[blue, very thick, samples=80, domain=-10.5:-5.10000000000000, smooth] plot (\x, {-4 * (\x - (8)) * (\x - (-3)) / ((\x - (-5)) * (\x - (4)))}); \draw[blue, very thick, samples=80, domain=-4.90000000000000:3.90000000000000, smooth] plot (\x, {-4 * (\x - (8)) * (\x - (-3)) / ((\x - (-5)) * (\x - (4)))}); \draw[blue, very thick, samples=80, domain=4.10000000000000:10.5, smooth] plot (\x, {-4 * (\x - (8)) * (\x - (-3)) / ((\x - (-5)) * (\x - (4)))}); \fill[blue] (8.0, 0) circle (4pt); \fill[blue] (-3.0, 0) circle (4pt); \fill[blue] (0, -4.8) circle (4pt); \fill[blue] (-9.0, -7.85) circle (4pt); \fill[blue] (-4.0, 6.0) circle (4pt); \fill[blue] (6.5, 1.98) circle (4pt); \draw[blue, very thick, fill=white] (-8.0, -8.89) circle (4pt); \end{tikzpicture}\] 
\end{solution}

    \vspace{\stretch{1}}

\question[6] Find the exact value of each of the following expressions.\\
\begin{multicols}{2}
\begin{parts}
    \part $\sin \left(-\dfrac{3\pi}{4}\right)$ = \hrulefill \vspace{12pt}
    \part $\sin \left(-\dfrac{3\pi}{4}\right)$ = \hrulefill \vspace{12pt}
    \part $\sin \left(-\dfrac{3\pi}{4}\right)$ = \hrulefill \vspace{12pt}
    \part $\cos (210^\circ)$ = \hrulefill \vspace{36pt}
    \part $\cos (210^\circ)$ = \hrulefill \vspace{36pt}
    \part $\cos (210^\circ)$ = \hrulefill 
    \end{parts}
\end{multicols}

\nocalc{\version}

\question[10] Find the exact value(s) of $\theta$ in the indicated units.
\begin{multicols}{2}
\begin{parts}
    \uplevel{\textbf{Degrees: $[0^\circ,360^\circ)$}} 
    \vspace{12pt}
    
    \part $\cot\theta=-1$, so $\theta=$ \hrulefill\\  
    \vspace{12pt}
    
    \part $\cos\theta=\dfrac{1}{2}$, so $\theta=$ \hrulefill\\ 
    \vspace{12pt}

    \part $\cos\theta=\dfrac{1}{2}$, so $\theta=$ \hrulefill\\ 
    \vspace{12pt}    

\columnbreak

    \uplevel{\textbf{Radians: $[0,2\pi)$}} 
    \vspace{12pt}
    
    \part $\cos\theta=0$, so $\theta=$ \hrulefill\\ 
    \vspace{12pt}
    
    \part $\sin\theta=-\dfrac{\sqrt{2}}{2}$, so $\theta=$ \hrulefill\\ 
    \vspace{12pt}
    
    \part $\sin\theta=-\dfrac{\sqrt{2}}{2}$, so $\theta=$ \hrulefill\\ 
    \vspace{12pt}\end{parts}
\end{multicols}

\question[2]
Suppose a trigonometric function has the attributes in the table given below. Using only the information provided, sketch the \textbf{transformed wave of interest} for the function of your choice. \textit{\textbf{Be careful not to include any changes to the function other than the ones listed.}}

\begin{multicols}{2}
    Select \textbf{one} of the following functions to sketch. \textbf{\textit{You MUST indicate your selection to earn credit}}.
    \begin{checkboxes}
    \choice Sine   
    \choice Cosine    
    \choice Tangent
\end{checkboxes}

\columnbreak
\hspace{1in}{
    \renewcommand{\arraystretch}{2}
\begin{tabular}{ |r|c| } 
 \hline
 Period: & $4$ \\\hline
 Vert. Dilation: & $2$\\\hline
 Phase Shift: & $-1$ \\\hline
\end{tabular}}

\end{multicols}

\vspace{0.25in}

\begin{center}
    \begin{tikzpicture}[scale=.75,>=triangle 45] 
        % Draw horizontal axis
        \draw[<->, very thick] (-10, 0) -- (10, 0);
    
        % Draw vertical axis
        \draw[<->, very thick] (0, -3) -- (0, 3);
    \end{tikzpicture}
\end{center}
\vspace{0.25in}

\nocalc{\version}

\question[3]
Given \(\sin \theta = \dfrac{4}{5}\) and with \(\dfrac{\pi}{2} < \theta < \pi\), use identities to find \(\sin(2\theta)\).\\
\begin{solution}
    \(\sin(2\theta) = 2\sin\theta\cos\theta\)
    \(\sin\theta = \frac{4}{5} \quad \cos\theta = -\frac{3}{5}\)
    \(= 2 \left(\frac{4}{5}\right) \left(-\frac{3}{5}\right)\)
    \(\boxed{-\frac{24}{25}}\)
\end{solution}
    \vspace{\stretch{1}}
    \answerline

\question[3]
Write the composition of trigonometric functions as an algebraic (non-trigonometric) expression in \(u\) with \(u > 0\).\\

\(\cos\left(2\tan^{-1}\left(3u\right)\right)\)
    \begin{solution}
        \(\text{Let } \theta = \tan^{-1}\left(3u\right) \implies \tan\theta = \frac{3u}{1} = \frac{y}{x}\)
        \(r = \sqrt{x^2 + y^2} = \sqrt{(1)^2 + (3u)^2} = \sqrt{9u^2 + 1}\)
        \(\cos\theta = \frac{x}{r} = \frac{1}{\sqrt{9u^2 + 1}}\)
        \(\cos(2\theta) = 2\cos^2\theta - 1\)
        \(= 2\left(\frac{1}{\sqrt{9u^2 + 1}}\right)^2 - 1\)
        \(= 2 \cdot \frac{1}{9u^2 + 1} - 1\)
        \(= \frac{2}{9u^2 + 1} - \frac{9u^2 + 1}{9u^2 + 1}\)
        \(= \frac{2 - 9u^2 - 1}{9u^2 + 1}\)
        \(\boxed{\frac{1 - 9u^2}{9u^2 + 1}}\)
    \end{solution}
    \vspace{\stretch{1}}
    \answerline

\nocalc{\version}

\question[3]
Solve the following trigonometric equation on the interval \([0, 2\pi)\).\\ \\
\(2\cos^2 \theta - 7\cos \theta = (-3)\)
    \begin{solution}
        \(2\cos^2 \theta - 7\cos \theta + 3 = 0\)
        \((2\cos \theta - 1)(\cos \theta - 3) = 0\)
        \(2\cos \theta - 1 = 0 \quad \text{OR} \quad \cos \theta - 3 = 0\)
        \(\cos \theta = \frac{1}{2} \quad \text{OR} \quad \cos \theta = 3\)
        \(\cos \theta = \frac{1}{2} \quad \text{OR} \quad \text{No real solution}\)
        \(\boxed{\theta = \frac{1}{3} \, \pi, \frac{5}{3} \, \pi}\)
    \end{solution}
    \vspace{\stretch{1}}
    \answerline

\question[3]
Solve for exact values of \(x\).\\

\(\cos^{-1}(x) - \tan^{-1}\left(-1\right) = \dfrac{3\pi}{4}\)
    \begin{solution}
        \(\cos^{-1}(x) = \frac{3\pi}{4} - \frac{\pi}{4}\)
        \(\cos^{-1}(x) = \frac{\pi}{2}\)
        \(x = \cos\left(\frac{\pi}{2}\right)\)
        \(\boxed{x = 0}\)
    \end{solution}
    
    \vspace{\stretch{1}}
    \answerline

\nocalc{\version}

\question[1] Plot \textbf{\textit{and label}} the given polar coordinates.
\begin{multicols}{2}
\begin{parts}
    \vspace*{.01 in}
    \part \(\left(-4, \dfrac{7\pi}{6}\right)\)\vspace{.5 in}
    \part \(\left(-2, 180^\circ\right)\)\vspace{.5 in}
    \vspace*{.5 in}
\end{parts}
\columnbreak
\includegraphics[width=3in]{10. Final/Images/Polar.pdf}
\end{multicols}

\question[2] Convert the previous polar coordinates to rectangular coordinates.
    \begin{parts}
    \part \(\left(-4, \dfrac{7\pi}{6}\right)\)\vspace{.5 in}
            \begin{solution}
                \(A: (x, y) = \left(2 \, \sqrt{3}, 2\right)\)
            \end{solution}
        \vspace{\stretch{1}}
        \answerline
    \part \(\left(-2, 180^\circ\right)\)\vspace{.5 in}
            \begin{solution}
                \(B: (x, y) = \left(2, 0\right)\)
            \end{solution}
        \vspace{\stretch{1}}
        \answerline\end{parts}

\nocalc{\version}

\uplevel{Reference the function $f(x)$ graphed below for questions \ref{graph-limit-start} through \ref{graph-limit-end}. If the limit or function value does not exist, write \emph{DNE}.}

\begin{center}
\begin{tikzpicture}[scale=0.6,>=triangle 45]
    % Grid
    \draw[style={black!40}] (-7,-5) grid (7,5);

    % Asymptotes
    \draw[<->, very thick, dashed, orange] (1,-5.5) -- (1,5.5);

    % Axes
    \draw[<->, very thick] (-7.5,0) -- (7.5,0);
    \draw[<->, very thick] (0,-5.5) -- (0,5.5);

    % Y-axis labels
    \foreach \y in {-4,-2,2,4} {
        \node[left=2pt, fill=white, inner sep=1pt] at (0,\y) {\scriptsize \y};}

    % X-axis labels
    \foreach \x in {-6,-4,-2,2,4,6} {
        \node[below=2pt, fill=white, inner sep=1pt] at (\x,0) {\scriptsize \x};}

    % Function plot pieces
    % Left parabola piece
    \draw[<-, ultra thick, style={blue}, domain=-6.4:-2, samples=50]
        plot (\x, {1.25*(\x+4)^2 - 2});
        
    % Middle linear segment
    \draw[-, ultra thick, style={blue}] (-2,-2) -- (0,0);
    
    % Left side of asymptote
    \draw[->, ultra thick, style={blue}, domain=0:0.833, samples=50]
        plot (\x, {1/(\x-1) + 1});
        
    % Right side of asymptote
    \draw[<-, ultra thick, style={blue}, domain=1.285:3, samples=50]
        plot (\x, {-2/(\x-1) + 2});
        
    % Right parabola piece
    \draw[->, ultra thick, style={blue}, domain=3:7.1, samples=50]
        plot (\x, {0.25*(\x-3)^2 + 1});
        
    % Points
    \filldraw[fill=white, draw=blue, very thick] (-2,3) circle (0.2);
    \filldraw[blue] (-2,-2) circle (0.2);
    \filldraw[blue] (0,0) circle (0.2);
    \filldraw[fill=white, draw=blue, very thick] (3,1) circle (0.2);
    % \filldraw[blue] (3,-2) circle (0.2);
\end{tikzpicture}\end{center}    

\question[2]\,\\
\begin{parts}   
\part
  $\displaystyle{\lim_{x \to (-2)^-} f(x)}=$ \\ \\  
   \vspace*{.15in}
\part
  $\displaystyle{\lim_{x \to (-2)^+} f(x)}=$ \\ \\  
   \vspace*{.15in}
\part
  $\displaystyle{\lim_{x \to (-2)} f(x)}=$ \\ \\  
   \vspace*{.15in}
\part
  $f(-2)=$ \\ \\
   \vspace*{.15in}
\end{parts}

\label{graph-limit-start}
  
\question[1]
  $\displaystyle{\lim_{x \to (-\infty)} f(x)}=$ \\ \\
   \vspace*{.15in}

\question[1]
  For what value $a$ does $f(a)$ not exist, but $\displaystyle{\lim_{x \to a} f(x)}$ does?
   \vspace*{.15in}

\label{graph-limit-end}

\end{questions}



%%%%%%%%%%%%%%%%%%%%%%%%%% Scratch Paper %%%%%%%%%%%%%%%%%%%%%%%%%%

\UCblankpage

\blankpage

\end{document}
